%% Preamble (outline.tex)
\usepackage{microtype} % Micro-typographic refinements

% Core math & theorem packages (unicode-math replaces amsfonts/amssymb)
\usepackage{amsmath}
\usepackage{amsthm}
\usepackage{extarrows}

% Tables, lists, graphics, colors
\usepackage{booktabs}
\usepackage{xcolor}
\usepackage{multirow}
\usepackage{multicol}
\usepackage{enumitem}
\usepackage{appendix}
\usepackage{longtable}
\usepackage{threeparttable}
\usepackage{tabularx}
\usepackage{graphicx}

% Fonts (XeLaTeX — Libertinus family, loaded directly from TeX Live OTF files)
\usepackage{fontspec}
\setmainfont{LibertinusSerif-Regular.otf}[
    BoldFont=LibertinusSerif-Bold.otf,
    ItalicFont=LibertinusSerif-Italic.otf,
    BoldItalicFont=LibertinusSerif-BoldItalic.otf
]
\setsansfont{LibertinusSans-Regular.otf}[
    BoldFont=LibertinusSans-Bold.otf,
    ItalicFont=LibertinusSans-Italic.otf
]
\setmonofont{LibertinusMono-Regular.otf}[Scale=0.88]
\usepackage{unicode-math}
\setmathfont{LibertinusMath-Regular.otf}

% Chinese support for Appendix E (Bilingual Table) under XeLaTeX
\usepackage{xeCJK}
\setCJKmainfont{FandolKai-Regular.otf}
\setCJKsansfont{FandolKai-Regular.otf}
\setCJKmonofont{FandolKai-Regular.otf}
\newfontfamily{\kai}{FandolKai-Regular.otf}

% Page geometry & headers
\usepackage{geometry}
\geometry{a4paper, left=2.2cm, right=2.2cm, top=2.4cm, bottom=2.4cm, headheight=15pt}

\usepackage{fancyhdr}
\pagestyle{fancy}
\fancyhf{}
\fancyhead[ER]{\leftmark}
\fancyhead[OL]{\rightmark}
\fancyhead[EL,OR]{\thepage}
\renewcommand{\headrulewidth}{0.4pt}
\renewcommand{\footrulewidth}{0pt}

% Section & ToC formatting
\usepackage{titlesec}
\titleformat{\chapter}{\Huge\bfseries}{\thechapter}{1em}{}
\titleformat{\section}{\Large\bfseries}{\thesection}{1em}{}
\titleformat{\subsection}{\large\bfseries}{\thesubsection}{1em}{}
\titleformat{\subsubsection}{\normalsize\bfseries}{\thesubsubsection}{1em}{}
\titleformat{\paragraph}{\normalsize\bfseries}{\theparagraph}{1em}{}

\usepackage{titletoc}
\titlecontents{chapter}[0em]{\bfseries}{\thecontentslabel\quad}{}{\titlerule*[10pt]{$\cdot$}\contentspage}
\titlecontents{section}[2em]{}{\thecontentslabel\quad}{}{\titlerule*[10pt]{$\cdot$}\contentspage}
\titlecontents{subsection}[4em]{}{\thecontentslabel\quad}{}{\titlerule*[10pt]{$\cdot$}\contentspage}

% Theorem box colors
\definecolor{theoremBG}{rgb}{0.8, 0.89, 1}          % #CCE2FF
\definecolor{theoremTitleBG}{rgb}{0.43, 0.7, 0.86}  % #6EB2DB
\definecolor{definitionBG}{rgb}{0.86, 0.95, 0.98}   % #DBF2F9
\definecolor{definitionTitleBG}{rgb}{0.55, 0.89, 0.89} % #8BE3E3
\definecolor{mistyrose}{rgb}{1, 0.89, 0.88}
\definecolor{seashell}{rgb}{1, 0.96, 0.93}
\definecolor{lemmaBG}{rgb}{0.8, 0.94, 0.96}         % #CBF1F5
\definecolor{lemmaTitleBG}{rgb}{0.65, 0.89, 0.91}   % #A6E3E9
\definecolor{corollaryBG}{rgb}{0.74, 0.79, 0.92}    % #BDCAEA
\definecolor{corollaryTitleBG}{rgb}{0.53, 0.56, 0.8} % #8890CC
\definecolor{lightblue}{rgb}{0.8, 0.9, 1}

% Framed theorem environments
\usepackage{mdframed}

\newmdtheoremenv[
    linecolor=definitionTitleBG,
    backgroundcolor=definitionBG,
    linewidth=3pt,
    topline=false,
    bottomline=false,
    rightline=false,
    roundcorner=0pt,
    frametitlebackgroundcolor=definitionTitleBG,
    frametitlefont=\bfseries,
    innertopmargin=2pt,
    innerbottommargin=8pt,
    innerrightmargin=8pt,
    innerleftmargin=8pt
]{definition}{Definition}[chapter]

\newmdtheoremenv[
    linecolor=mistyrose,
    backgroundcolor=seashell,
    linewidth=3pt,
    topline=false,
    bottomline=false,
    rightline=false,
    roundcorner=0pt,
    frametitlebackgroundcolor=mistyrose,
    frametitlefont=\bfseries,
    innertopmargin=2pt,
    innerbottommargin=8pt,
    innerrightmargin=8pt,
    innerleftmargin=8pt
]{proposition}{Proposition}[chapter]

\newmdtheoremenv[
    linecolor=theoremTitleBG,
    backgroundcolor=theoremBG,
    linewidth=3pt,
    topline=false,
    bottomline=false,
    rightline=false,
    roundcorner=0pt,
    frametitlebackgroundcolor=theoremTitleBG,
    frametitlefont=\bfseries,
    innertopmargin=2pt,
    innerbottommargin=8pt,
    innerrightmargin=8pt,
    innerleftmargin=8pt
]{theorem}{Theorem}[chapter]

\newmdtheoremenv[
    linecolor=lemmaTitleBG,
    backgroundcolor=lemmaBG,
    linewidth=3pt,
    topline=false,
    bottomline=false,
    rightline=false,
    roundcorner=0pt,
    frametitlebackgroundcolor=lemmaTitleBG,
    frametitlefont=\bfseries,
    innertopmargin=2pt,
    innerbottommargin=8pt,
    innerrightmargin=8pt,
    innerleftmargin=8pt
]{lemma}{Lemma}[chapter]

\newmdtheoremenv[
    linecolor=corollaryTitleBG,
    backgroundcolor=corollaryBG,
    linewidth=3pt,
    topline=false,
    bottomline=false,
    rightline=false,
    roundcorner=0pt,
    frametitlebackgroundcolor=corollaryTitleBG,
    frametitlefont=\bfseries,
    innertopmargin=2pt,
    innerbottommargin=8pt,
    innerrightmargin=8pt,
    innerleftmargin=8pt
]{corollary}{Corollary}[chapter]

% Upright theorem environments numbered by chapter
\theoremstyle{definition}
\newtheorem{example}{Example}[chapter]
\newtheorem{note}{Note}[chapter]
\newtheorem{remark}{Remark}[chapter]
\newtheorem{problem}{Problem}[chapter]

% Unified mathematical macros
\newcommand{\E}{\mathbb{E}}
\newcommand{\Prob}{\mathbb{P}}
\newcommand{\Q}{\mathbb{Q}}
\newcommand{\R}{\mathbb{R}}
\newcommand{\N}{\mathbb{N}}
\newcommand{\F}{\mathcal{F}}
\newcommand{\dd}{\,\mathrm{d}}
\newcommand{\ind}{\mathbb{1}}
\DeclareMathOperator{\Var}{Var}
\DeclareMathOperator{\Cov}{Cov}
\DeclareMathOperator{\BS}{BS}

% Load hyperref last
\usepackage[hidelinks]{hyperref}
